\section{Revisão de Literatura}\label{sec:revisao}

Esta seção reúne o necessário para compreender a solução: como a análise de solo é
interpretada, como o Manual de Calagem e Adubação transforma a interpretação em recomendação,
como se avalia a aptidão das terras e o que caracteriza um sistema especialista baseado em
regras. Ao final, posiciona-se o trabalho frente às ferramentas existentes.

\subsection{Análise de solo e sua interpretação}\label{sec:rev_analise}

A análise de solo é o principal instrumento de diagnóstico da fertilidade. Nela se determinam
o pH em água, que condiciona a disponibilidade da maioria dos nutrientes; os teores de fósforo
(P) e potássio (K) disponíveis; os teores trocáveis de cálcio (Ca), magnésio (Mg) e alumínio
(Al), este associado à toxidez em solos ácidos; a CTC, que expressa a capacidade do solo de
reter cátions; a saturação por bases (V\%); o teor de matéria orgânica (MO), relacionado ao
suprimento de nitrogênio; e o teor de argila, principal atributo físico considerado na
recomendação \citep{Donagema2011, ManualCalagem2016}.

Os valores, isoladamente, não dizem nada: a interpretação os compara a faixas de referência
que dependem do tipo de solo e da cultura \citep{EmbrapaDoc206}. No Manual RS/SC, os teores de
P e K são enquadrados em cinco classes de disponibilidade --- muito baixo, baixo, médio, alto
e muito alto ---, e cada classe corresponde a uma faixa de rendimento relativo da cultura; o
limite superior da classe médio é o teor crítico, acima do qual se obtém rendimento próximo
ao máximo econômico \citep{ManualCalagem2016}. As faixas não são universais: a interpretação
do fósforo depende da classe de teor de argila do solo, e a do potássio, da CTC a pH 7,0. Um
mesmo teor de P pode, assim, ser baixo em um solo argiloso e alto em um arenoso.

\subsection{Calagem e adubação no RS e em SC}\label{sec:rev_manual}

O Manual de Calagem e Adubação para os estados do Rio Grande do Sul e de Santa Catarina
\citep{ManualCalagem2016} é a referência técnica para correção e fertilização do solo na
região. A necessidade de calcário é estimada preferencialmente pelo \textbf{índice SMP}, que
indica a dose de calcário com PRNT 100\% necessária para elevar o pH da camada de 0 a 20~cm a
5,5, 6,0 ou 6,5; o método da saturação por bases é previsto como alternativa. Dois pontos
tornam essa recomendação menos direta do que uma consulta a tabela. Primeiro, o pH que dispara
a calagem e o pH alvo da dose não são dados pelo pH de referência da cultura, e frequentemente
divergem dele: são definidos por grupo de culturas e por sistema de manejo --- grãos em sistema convencional, por exemplo, só recebem
calcário quando o pH é inferior a 5,5, com dose para pH 6,0. Segundo, a dose é ajustada ao
contexto: em plantio direto consolidado aplica-se um quarto da dose, em superfície e limitada
a 5~t/ha; solos de baixo poder tampão usam equações com MO e Al em lugar da tabela; e a dose
é corrigida pelo PRNT do corretivo efetivamente disponível \citep{ManualCalagem2016}.

A adubação fosfatada e potássica é dimensionada pela classe de disponibilidade e, para várias
culturas, pela expectativa de rendimento; a nitrogenada, pelo teor de MO e pela cultura
antecedente, e é dispensada nas leguminosas, que fixam nitrogênio biologicamente. O Manual
publica tabelas próprias para cada cultura ou grupo, com estruturas distintas: grãos separam
adubação de correção e de manutenção, frutíferas distinguem pré-plantio, crescimento e
produção, e a erva-mate e o tabaco exigem variáveis específicas de manejo e de tipo
\citep{ManualCalagem2016}.

\subsection{Aptidão das terras e o recorte edáfico}\label{sec:rev_aptidao}

O Sistema de Avaliação da Aptidão Agrícola das Terras \citep{RamalhoFilhoBeek1995} classifica
as terras pelo grau de limitação --- nulo, ligeiro, moderado, forte e muito forte --- que cada
fator impõe, considerando deficiência de fertilidade, deficiência hídrica, excesso de água,
suscetibilidade à erosão e impedimentos à mecanização. A classe resulta do fator mais
limitante, avaliado nas condições atuais e após a aplicação de práticas de melhoria
\citep{Hofig2015}. É o sistema empregado no mapeamento nacional \citep{EmbrapaAptidao2025}.
Sua escala, porém, é regional: \cite{Hofig2015} observam que ele não indica práticas de manejo
e requer adaptações para o planejamento de propriedades individuais.

Uma alternativa é restringir a avaliação aos atributos que a própria análise de solo informa,
cruzados com as exigências da cultura, sem relevo e sem clima. Neste trabalho, essa estimativa
é chamada de \textbf{aptidão edáfica}, para que não seja confundida com a aptidão agrícola das
terras: ela herda a arquitetura de graus de limitação e de fator mais limitante, mas não
reproduz o sistema oficial nem equivale a ele.

\subsection{Sistemas especialistas baseados em regras}\label{sec:rev_se}

Sistemas especialistas reproduzem o raciocínio de um especialista humano em um domínio
restrito \citep{Mendes1997}. São compostos por uma \textbf{base de conhecimento}, que
armazena fatos e regras do domínio; um \textbf{motor de inferência}, que aplica as regras aos
dados de entrada para derivar conclusões; e uma \textbf{interface} com o usuário. Na forma mais
difundida, o conhecimento é representado por regras SE-ENTÃO, e o motor percorre a base para
identificar as condições satisfeitas pelos dados \citep{Lanzer2009}. A representação explícita
das regras é o que permite ao sistema \emph{explicar} uma conclusão --- característica central
em domínios nos quais a recomendação precisa ser justificada a terceiros.

Na agricultura, \cite{Lanzer2009} aplicou representação hierárquica do conhecimento à
classificação do uso e manejo das terras e observou que, quando os estados do domínio já estão
definidos em sistemas de classificação publicados --- como as tabelas do Manual ---, a
construção da base é facilitada. \cite{Shahzadi2016}, em outro domínio, alimentaram regras
especialistas com dados de sensores para recomendar o manejo de pragas e doenças, evidenciando a
aplicabilidade da abordagem ao suporte à decisão agronômica.

\subsection{Trabalhos relacionados}\label{sec:rev_relacionados}

O \textbf{FertFácil} \citep{Tiecher2016} é um sistema on-line e gratuito que recomenda calagem
e adubação pelos critérios do Manual para mais de cem cultivos do RS, SC, MS e Paraguai, adotado
pela Emater/RS; hoje incorpora também o manual do Paraná \citep{FertFacil2026}. O \textbf{AdubaTec} \citep{AdubaTec2021}, da Embrapa Mandioca e Fruticultura e da
UFRB, recomenda calagem e adubação para mandioca e fruteiras tropicais a partir das recomendações
publicadas pela própria Embrapa, com arquitetura parametrizável que permite incluir novas culturas
sem alterar o código. O \textbf{Sistema de Avaliação da Aptidão
Agrícola das Terras} \citep{RamalhoFilhoBeek1995, EmbrapaAptidao2025} classifica o potencial
das terras em escala regional. A \autoref{tab:relacionados} sintetiza o posicionamento: os
sistemas existentes tratam separadamente a recomendação de manejo e a avaliação de aptidão, e
nenhum integra as duas a partir de uma análise individual de solo sob os critérios do RS e SC.

\begin{table}[htb]
\centering
\caption{Posicionamento do SIRAS frente aos trabalhos relacionados.}
\label{tab:relacionados}
\footnotesize
\begin{tabularx}{\columnwidth}{@{}lCCCC@{}}
\toprule
 & \textbf{Critérios} & \textbf{Escala} & \textbf{Calagem e adubação} & \textbf{Aptidão} \\
\midrule
FertFácil  & RS/SC e PR & área & sim & não \\
AdubaTec   & Embrapa & área  & sim & não \\
SAAAT      & nacional & regional & não & sim \\
\textbf{SIRAS} & \textbf{RS/SC} & \textbf{área} & \textbf{sim} & \textbf{sim (edáfica)} \\
\bottomrule
\end{tabularx}

\vspace{2pt}
\raggedright\scriptsize SAAAT: Sistema de Avaliação da Aptidão Agrícola das Terras.
\end{table}
